\documentclass[a4paper,10pt]{article}
\usepackage[margin=0.5in,nofoot]{geometry}
\usepackage{fontawesome5}
\usepackage{hyperref}
\usepackage{titlesec}
\usepackage{xcolor}
\usepackage{needspace}

\hypersetup{
    colorlinks=true,
    linkcolor=black,
    filecolor=black,
    urlcolor=black,
    citecolor=black
}

\titleformat{\section}{\large\bfseries}{\thesection}{1em}{}[\titlerule]
\titlespacing*{\section}{0pt}{*1}{*1}

\newcommand{\entry}[4]{
  \noindent\textbf{#1} \hfill #2 \\
  \noindent\textit{#3} \hfill \textit{#4} \\
  \vspace{2pt}
}

\newcommand{\project}[2]{
  \noindent\textbf{#1} \hfill #2 \\
  \vspace{2pt}
}

\begin{document}

\pagenumbering{gobble}

\noindent
\begin{minipage}[t]{0.5\textwidth}
\textbf{\Large Danilo Chagas Clemente}

\vspace{0.4em}

\end{minipage}%
\begin{minipage}[t]{0.5\textwidth}
\raggedleft

Lavras - MG, Brasil

{\color{blue}} \href{tel:+5535999840288}{\faPhone \space (35) 99984-0288}
{\color{blue}} \href{mailto:danilo.chagasclem@gmail.com}{\faEnvelope \space danilo.chagasclem@gmail.com}

\vspace{0.2em}
 \quad
{\color{blue}} \href{https://github.com/danilochagasc}{\faGithub \space GitHub} \quad
{\color{blue}} \href{https://www.linkedin.com/in/danilo-clemente-94880a232/}{\faLinkedin \space LinkedIn} \\
\end{minipage}

\vspace{1em}

\noindent\textbf{\Large Engenheiro de Software}
\vspace{0.5em}

\section*{Resumo Profissional}

\vspace{0.6em}

Engenheiro de Software Backend com experiência em sistemas de alta escala, processando milhões de transações diárias. Bacharel em Ciência da Computação pela UFLA e mestrando na mesma instituição. Domínio em Java, Kotlin (Spring Boot), TypeScript (Angular) e Python, com vivência em arquiteturas de microsserviços, mensageria (Kafka, SQS), bancos SQL/NoSQL e práticas de observabilidade. Forte capacidade analítica, comunicação clara e foco em soluções robustas e escaláveis.

\vspace{0.6em}

\section*{Experiência Profissional}

\vspace{0.6em}

\entry{iFood}{\faCalendar \space Jan/2026 -- Atual}{Engenheiro de Software Backend}{\faMapMarker \space Remoto}
\space
\vspace{-1.6em}
\begin{itemize}
\setlength\itemsep{0em}
\item Desenvolvo e mantenho microsserviços em Kotlin e Spring Boot na plataforma de pagamentos, responsável por serviços que processam mais de 6 milhões de transações por dia.
\item Integrei as plataformas iFood e Uber, contribuindo para a interoperabilidade entre sistemas e expansão de funcionalidades para milhões de usuários.
\item Implemento e mantenho serviços financeiros, incluindo meios de pagamento, conciliação e integrações com adquirentes.
\item Desenvolvo testes unitários e de integração com JUnit e Kotest, garantindo qualidade e confiabilidade em ambientes de alta disponibilidade.
\item Monitoro e resolvo incidentes em produção com Datadog, analisando métricas, logs e traces para manter SLAs de disponibilidade.
\item Projeto soluções com comunicação assíncrona (Kafka, SQS) e arquiteturas distribuídas, aplicando arquitetura hexagonal, DDD e CQRS.
\item Gerencio containerização com Docker e orquestração com Kubernetes em ambiente AWS.
\item Executo deploys com estratégias de canary release, reduzindo riscos e garantindo rollback seguro em produção.
\end{itemize}

\entry{iFood}{\faCalendar \space Fev/2025 -- Dez/2025}{Estagiário de Engenharia de Software}{\faMapMarker \space Remoto}
\space
\vspace{-1.6em}
\begin{itemize}
\setlength\itemsep{0em}
\item Desenvolvi e mantive APIs backend em Kotlin e Spring Boot em sistemas de pagamentos de alta escala.
\item Automatizei a retokenização de cartões com falha via script Python executado por cron job, aumentando a taxa de recuperação de cartões e eficiência operacional.
\item Configurei e dei suporte a meios de pagamento e integrações com adquirentes.
\item Implementei testes automatizados e apoiei o monitoramento de serviços em produção com Datadog.
\item Atuei com comunicação assíncrona (Kafka) e sistemas distribuídos em arquitetura de microsserviços.
\end{itemize}

\entry{Emakers Jr}{\faCalendar \space Jun/2023 -- Abr/2025}{Desenvolvedor Fullstack}{\faMapMarker \space Lavras - MG}
\space
\vspace{-1.6em}
\begin{itemize}
\setlength\itemsep{0em}
\item Desenvolvi funcionalidades fullstack com React (TypeScript) no frontend e Java (Spring Boot) no backend do sistema Redação Inteligente, plataforma de gestão de cursos e correção de redações com milhares de usuários ativos.
\item Migrei a integração de pagamentos de PagBank para Mercado Pago, incluindo suporte a Pix, aumentando a confiabilidade e modernizando o sistema de cobranças.
\item Conduzi a evolução da arquitetura backend, migrando de AdonisJS para Java com Spring Boot e de MySQL para PostgreSQL.
\item Desenvolvi soluções internas para plataformização de processos de gestão de pessoas e recrutamento, otimizando a eficiência operacional da organização.
\item Criei materiais técnicos e treinamentos (documentação em Markdown e cursos em vídeo) para capacitação interna em tecnologias backend.
\end{itemize}

\Needspace{14\baselineskip}
\entry{Nexos Digital}{\faCalendar \space Out/2024 -- Jan/2025}{Desenvolvedor Fullstack (Estágio)}{\faMapMarker \space Remoto}
\space
\vspace{-1.6em}
\begin{itemize}
\setlength\itemsep{0em}
\item Desenvolvi funcionalidades fullstack com Java (Spring Boot) e Angular (TypeScript) na plataforma Mercap Solutions para distribuição de produtos de renda fixa.
\item Implementei automação de mailing com atualizações de mercado, aumentando o engajamento de usuários da plataforma.
\item Desenvolvi módulos para gestão de dados e disponibilização de vídeos institucionais de ofertas.
\item Criei funcionalidades para geração e download de planilhas automatizadas, otimizando o processo de cadastro e gestão de ofertas.
\end{itemize}

\section*{Educação}
\vspace{0.6em}

\entry{Universidade Federal de Lavras (UFLA)}{\faCalendar \space 2026 -- Atual}{Mestrado em Ciência da Computação}{\faMapMarker \space Lavras - MG}

\entry{Universidade Federal de Lavras (UFLA)}{\faCalendar \space 2022 -- 2025}{Bacharelado em Ciência da Computação}{\faMapMarker \space Lavras - MG}

\entry{Centro Federal de Educação Tecnológica de Minas Gerais (CEFET-MG)}{\faCalendar \space 2019 -- 2021}{Técnico em Informática}{\faMapMarker \space Varginha - MG}

\section*{Habilidades Técnicas}
\vspace{0.6em}
\begin{itemize}
\setlength\itemsep{0em}
\item \textbf{Linguagens:} Java, Kotlin, Python, TypeScript, C/C++.
\item \textbf{Frameworks:} Spring Boot, Angular, React.
\item \textbf{Arquitetura:} Microsserviços, Arquitetura Hexagonal, DDD, CQRS, REST APIs, sistemas distribuídos.
\item \textbf{Mensageria:} Apache Kafka, AWS SQS/SNS.
\item \textbf{Banco de Dados:} PostgreSQL, MySQL, Redis, bancos NoSQL.
\item \textbf{Testes:} JUnit, Mockito, Kotest, testes unitários e de integração.
\item \textbf{DevOps e Cloud:} Docker, Kubernetes, AWS, CI/CD, canary release, Linux.
\item \textbf{Observabilidade:} Datadog, Grafana, Prometheus, análise de logs, métricas e traces.
\item \textbf{Ferramentas:} Git, GitHub, GitLab, IntelliJ IDEA.
\item \textbf{Processamento de Dados:} Pandas, Apache Spark.
\item \textbf{Metodologias:} Scrum, Kanban.
\item \textbf{Idiomas:} Português (nativo), Inglês (avançado).
\end{itemize}


\end{document}